% ==============================================================================
% SEZIONE 4: ANALISI DELLA PERFORMANCE ECONOMICO-FINANZIARIA
% ==============================================================================

\section{Il Caso Enel}

\begin{frame}{Analisi Finanziaria Enel: KPI FY 2025}
    La resilience finanziaria di Enel nel FY 2025 è il risultato di una gestione rigorosa del capitale e di un miglioramento qualitativo del mix dei ricavi.
    \vspace{0.1cm}
    
    \note{La performance economico finanziaria del Gruppo Enel nel 2025 evidenzia stabilità strutturale e rigida disciplina del capitale.}

    \begin{table}
        \centering
        \small
        \begin{tabular}{lc}
            \toprule
            \textbf{Metrica Finanziaria (FY 2025)} & \textbf{Valore / Impatto} \\
            \midrule
            \textbf{EBITDA Ordinario} & €22,9 mld (+2\% pro-forma) \\
            \textbf{Net Ordinary Income} & €7,0 mld (+6\%) \\
            \textbf{Indebitamento Finanziario Netto} & €57,2 mld \\
            \textbf{Leva Finanziaria (Net Debt / EBITDA)} & 2.5x \\
            \bottomrule
        \end{tabular}
    \end{table}
    \vspace{-0.1cm}
    \pause
    \note{L'EBITDA ordinario cresce del due per cento pro-forma, beneficiando del consolidamento delle attività retail e di rete.}

    \begin{block}{Solidità Patrimoniale}
        Il \textit{Net Ordinary Income} sale a €7,0 miliardi (+6%). La leva finanziaria si stabilizza a \textbf{2.5x}, garantendo piena capacità di autofinanziamento degli investimenti.
    \end{block}
    \pause
    \note{Il net ordinary income sale a sette miliardi di euro, mentre il rapporto debito netto su EBITDA si attesta a due virgola cinque volte.}

    \begin{block}{Focus Strategico Italia: Incremento della RAB}
        La crescita del 10\% in Italia si riferisce alla \textbf{Regulatory Asset Base (RAB)}, salita a \textbf{€23,2 miliardi}. Questo è il vero motore della visibilità dei margini futuri.
    \end{block}
    \note{In Italia la RAB cresce a ventitré miliardi e duecento milioni, consolidando la prevedibilità tariffaria del Gruppo.}
\end{frame}


\begin{frame}{Qualità dei Margini e Derisking Strutturale}
    Il riposizionamento strategico del Gruppo si riflette in una drastica riduzione dell'esposizione alla volatilità dei mercati commodity:
    \vspace{0.2cm}

    \begin{itemize}
        \item \textbf{Abbattimento del Rischio Merchant:} Il peso di Trading e \textit{Flexible Generation} sull'EBITDA è passato dal \textbf{\textasciitilde 50\% del 2022} al \textbf{\textasciitilde 20\% del 2025}.
        \pause
        \note{Il de risking strutturale riduce l'apporto di trading e generazione esposta al rischio merchant dal cinquanta al venti per cento.}
        
        \item \textbf{Efficienza Operativa e Qualità del Credito:} Gestione ottimale della voce \textit{D\&A e Bad Debt} (-€8,1 miliardi), con un netto miglioramento della qualità del credito.
        \pause
        \note{L'efficienza operativa ottimizza gli ammortamenti e svaluta i crediti per meno otto miliardi e cento milioni.}
        
        \item \textbf{Fidelizzazione Digitale:} Riduzione del \textbf{churn rate} di mercato grazie ai servizi avanzati e all'infrastruttura tecnologica.
    \end{itemize}
    \note{L'infrastruttura di rete e i servizi avanzati riducono il churn rate, consolidando la clientela retail. Passiamo alle infrastrutture di rete.}
\end{frame}
